\documentclass[11pt]{article}

\usepackage[margin=1in]{geometry}
\usepackage{amsmath,amssymb,amsthm}
\usepackage{mathtools}
\usepackage{enumitem}
\usepackage[colorlinks=true,linkcolor=blue,citecolor=blue,urlcolor=blue]{hyperref}

\newtheorem{theorem}{Theorem}
\newtheorem{lemma}[theorem]{Lemma}
\newtheorem{corollary}[theorem]{Corollary}
\newtheorem{proposition}[theorem]{Proposition}
\theoremstyle{definition}
\newtheorem{remark}[theorem]{Remark}
\newtheorem{definition}[theorem]{Definition}

\DeclareMathOperator{\ord}{ord}
\DeclareMathOperator{\lcm}{lcm}
\newcommand{\Fp}{\mathbb{F}_p}
\newcommand{\Z}{\mathbb{Z}}
\newcommand{\Q}{\mathbb{Q}}
\newcommand{\bin}[2]{\binom{#1}{#2}}
\newcommand{\wtd}{\widetilde}
\newcommand{\lean}[1]{\texttt{#1}}

\title{A machine-verified erratum and its repair:\\
Brown--Zudilin's $\zeta(5)$ denominators and the\\
Frobenius certificate, formalized in Lean~4}
\author{}
\date{19 July 2026. Revised: 20 July 2026.}

\begin{document}
\maketitle

\begin{abstract}
Brown and Zudilin's paper on cellular rational approximations to $\zeta(5)$
\cite{BZ22} states, as an experimental observation, the integrality
$d_n^5 P_n \in \Z$ for the rational coefficients of its totally symmetric
linear forms, where $d_n = \lcm(1,\dots,n)$. The paper's own displayed value
$P_2 = 1190161/384$ contradicts this inclusion, failing by a factor of exactly
$12$; the claim is still present, unamended, in the January 2026 revision. We
present a Lean~4 formalization, mechanically verified by the Lean kernel with
zero \lean{sorry}s and only the three standard axioms, of: (i)~the
contradiction, with $P_2$ \emph{derived} from the construction rather than
transcribed; (ii)~its sharp repair at $n = 2$: $c \, d_2^5 P_2 \in \Z$ if and
only if $12 \mid c$ (companion form: if and only if $2 \mid c$); and (iii)~the
theorem that proves the corrected law on the clean interval --- the
\emph{Frobenius certificate} theorem of \cite{CBcert}: for every $n \ge 1$
and every prime $p$ with $5 \le p$, $n < p \le 2n$, the harmonic functionals
satisfy $w_n \equiv \wtd{w}_n \equiv 0 \pmod p$, and hence
$p \mid p_n$. The formalization exposed two statement-level errors in our own
paper-level drafts, which we document. A second wave
(Section~\ref{sec:wave2}) adds: the corrected law kernel-verified for
$n \le 8$ with nonvanishing; a Lucas--Frobenius congruence
$Q_{ap+r} \equiv Q_a Q_r \pmod p$ for the cellular coefficient (whose
formalization caught a gap in the original proof sketch); and the
\emph{conditional assembly}: the descent hypothesis $(\mathrm D)$ implies the
full law at every prime $p \ge 5$ --- so the complete corrected law is now,
formally, one explicit hypothesis away. The complete development builds in
one \lean{lake build} and is publicly available.
\end{abstract}

\section{Introduction}

The irrationality theory of odd zeta values rests on rational linear forms in
zeta values whose coefficients must be cleared of denominators at a controlled
exponential cost: for a weight-$5$ form the benchmark cost is $d_n^5$, with
$d_n = \lcm(1,2,\dots,n) \sim e^n$. In the totally symmetric cellular
construction of Brown and Zudilin \cite{BZ22}, the linear form
\begin{equation}\label{eq:In}
I_n \;=\; Q_n\bigl(2\zeta(5) + 4\zeta(3)\zeta(2)\bigr)
\;-\; 4\widehat{P}_n\,\zeta(2) \;-\; 2P_n
\end{equation}
carries rational coefficient ladders $Q_n$, $\widehat{P}_n$, $P_n$, and the
paper states (in the display labeled \lean{dn-totsym}, presented explicitly as
an experimental observation from ``an extensive computation of the rational
coefficients''):
\begin{equation}\label{eq:claim}
Q_n,\;\; d_n^2 d_{2n}\widehat{P}_n,\;\; d_n^5 P_n \;\in\; \Z
\qquad (n = 0, 1, 2, \dots).
\end{equation}
The same paper displays (in the block following the definition of $I_n$) the
exact values
\begin{equation}\label{eq:displays}
Q_2 = 2989, \qquad
\widehat{P}_2 = \frac{344923}{96}, \qquad
P_2 = \frac{1190161}{384}.
\end{equation}
Since $d_2 = 2$ and $384 = 2^7\cdot 3$, we have
$d_2^5 P_2 = 1190161/12 \notin \Z$: the third inclusion of \eqref{eq:claim}
fails already at $n = 2$, by a factor of exactly $12$
(and the second fails by a factor of exactly $2$:
$d_2^2 d_4 \widehat{P}_2 = 344923/2$). The corrected law
\begin{equation}\label{eq:corrected}
12\, d_n^5 P_n \in \Z,
\end{equation}
conjecturally with the $\{2,3\}$-ceilings sharp, has substantial computational
support, and its clean-interval part is a theorem: the \emph{Frobenius
certificate} theorem of the companion paper \cite{CBcert} shows that every
prime $n < p \le 2n$ (the primes of $\bin{2n}{n}$, which the na\"ive
partial-fraction bound admits into the denominator) in fact divides the
numerator $p_n$.

This note reports a complete formalization of the erratum, its sharp $n=2$
repair, and the clean-interval theorem in Lean~4 over Mathlib. Every theorem
below is verified by the Lean kernel; the development contains zero
\lean{sorry}s, uses no \lean{native\_decide}, adds no axioms, and each named
theorem reports exactly the standard axiom triple
\lean{[propext, Classical.choice, Quot.sound]}. The trust base is the Lean
kernel and the printed displays of \cite{BZ22} quoted above --- nothing else.

Two purposes motivate machine verification here. First, an erratum to a
published paper should meet a higher evidentiary standard than the claim it
corrects; a kernel-checked derivation of $P_2$ from the construction, ending
in a formal negation of \eqref{eq:claim} at $n=2$, is the strongest form of
such a statement we know how to produce. Second, the formalization is a test
of a design thesis: proofs shaped as \emph{finite certificates plus one small
uniform identity} formalize at low cost. The clean-interval theorem has
exactly this shape, and the entire development --- definitions, numeric
gates, five lemma modules, assembly --- was completed in a single day. As a
byproduct the formalization caught two statement-level errors in our own
drafts (Section~\ref{sec:caught}), neither of which affects the mathematics,
both of which would likely have survived human refereeing.

\section{The construction and the formal objects}\label{sec:objects}

Fix $n \ge 1$ and set
\begin{equation}\label{eq:Rn}
R_n(k) \;=\; (n!)^4\,\Bigl(k+\frac n2\Bigr)\,
\frac{\prod_{j=1}^{n}(k-j)\,\prod_{j=1}^{n}(k+n+j)}{\prod_{j=0}^{n}(k+j)^6},
\end{equation}
Zudilin's totally symmetric rational function \cite{Zu02}. Its partial
fractions
\[
R_n(k) \;=\; \sum_{i=1}^{6}\sum_{j=0}^{n} \frac{a_{i,j}}{(k+j)^i},
\qquad
\wtd a_{i,j} \;=\; j(n-j)\,a_{i,j} + (2j-n)\,a_{i+1,j} - a_{i+2,j}
\]
(indices above $6$ read as zero; the $\wtd a$ are the partial fractions of
$-k(k+n)R_n(k)$) define the functionals
\[
u = \sum_j a_{5,j}, \quad
w = \sum_j a_{3,j}, \quad
v = \sum_{i,j} a_{i,j} H_j^{(i)}, \qquad
H_j^{(i)} = \sum_{m=1}^{j} m^{-i},
\]
their tilded companions, and the determinant ladders
\begin{equation}\label{eq:ladders}
q_n = u\wtd w - \wtd u w, \qquad
p_n = \wtd w v - w \wtd v, \qquad
\wtd p_n = u \wtd v - \wtd u v.
\end{equation}
The Brown--Zudilin coefficients of \eqref{eq:In} correspond via
$P_n = (-1)^{n+1} p_n / \bin{2n}{n}$ (the display following
\eqref{eq:claim} in \cite{BZ22}, which ties the cellular normalization to
\cite{Zu02}).

In Lean, all of these are \emph{computable} functions on $\Q$: the primary
definition of $a_{i,j}$ is the truncated local Taylor expansion of
$B_j(k) = (k+j)^6 R_n(k)$ at $k = -j$ (module \lean{Cbcert/Defs.lean}), and
the partial-fraction property is then a \emph{theorem}
(\lean{pf\_cleared}, module \lean{Cbcert/PartialFraction.lean}): the cleared
polynomial identity
\[
N_n(k) \;=\; \sum_{i,j} a_{i,j}\,(k+j)^{6-i}\prod_{m \ne j}(k+m)^6,
\]
proven for all $n$. Computability is what makes the development
self-auditing: a numeric gate (\lean{Cbcert/Numeric.lean}) reduces the full
definition chain by kernel computation at $n = 2$,
\[
w = \tfrac{6125}{4},\quad \wtd w = 1764,\quad u = 469,\quad \wtd u = 552,\quad
v = \tfrac{74463}{32},\quad \wtd v = \tfrac{43085}{16},\quad
p_2 = \wtd w v - w\wtd v = -\tfrac{1190161}{64},
\]
and at $n = 3$, in exact agreement with the independent reference
implementation; the frozen definitions were gated on this agreement before
any lemma work began. Note $P_2 = (-1)^{3} p_2/\bin{4}{2} = 1190161/384$,
matching \eqref{eq:displays} exactly.

\section{The machine-verified erratum}\label{sec:erratum}

\begin{theorem}[{\lean{Cbcert/ErrorExhibit.lean}}]\label{thm:erratum}
Define $P_2$ by the construction of Section~\ref{sec:objects}
\textup{(}derived, not transcribed\textup{)}. Then, verified by the Lean
kernel with standard axioms:
\begin{enumerate}[label=\textup{(\roman*)}]
\item \lean{P2\_eq}: $P_2 = 1190161/384$, and likewise
$Q_2 = 2989$, $\widehat P_2 = 344923/96$ --- the construction reproduces
\emph{all three} displayed values \eqref{eq:displays} of \cite{BZ22}.
\item \lean{bz\_inclusion\_fails}: $d_2^5\,P_2 \notin \Z$. In particular the
inclusion $d_n^5 P_n \in \Z$ of \eqref{eq:claim} is false at $n = 2$.
\item \lean{clearing\_iff}: for every $c \in \mathbb{N}$,
$\;c\,d_2^5\,P_2 \in \Z \iff 12 \mid c$. The repair factor $12$ of
\eqref{eq:corrected} is sharp at $n=2$.
\item \lean{companion\_inclusion\_fails}, \lean{companion\_clearing\_iff}:
$d_2^2 d_4 \widehat P_2 \notin \Z$, and
$c\,d_2^2 d_4 \widehat P_2 \in \Z \iff 2 \mid c$.
\end{enumerate}
\end{theorem}

The Lean statements are short enough to read in full (unicode transliterated):

\begin{quote}\small\ttfamily
theorem bz\_inclusion\_fails :\\
\hspace*{1em}$\lnot$ $\exists$ m : $\Z$, (d 2 : $\Q$)\^{}5 * P$_2$ = (m : $\Q$)\\[2pt]
theorem clearing\_iff (c : $\mathbb{N}$) :\\
\hspace*{1em}($\exists$ m : $\Z$, (c : $\Q$) * (d 2)\^{}5 * P$_2$ = m) $\leftrightarrow$ 12 $\mid$ c
\end{quote}

\noindent
where \lean{d} is the genuine running $\lcm(1,\dots,n)$. Each definition and
theorem carries a docstring citing the source labels of \cite{BZ22}
(\lean{dn-totsym} for the claim, the \lean{I\_n} display block for the
values), so a reader can hold the PDF next to the Lean file. Two remarks on
scope. The claim \eqref{eq:claim} is presented in \cite{BZ22} as an
experimental observation, not a theorem; what is exhibited here is a false
printed claim and its sharp correction, not a broken proof. And the bridge
between the cellular normalization of \cite{BZ22} and the construction of
Section~\ref{sec:objects} is the paper's own stated correspondence to
\cite{Zu02}; it is corroborated \emph{inside} the formalization by
Theorem~\ref{thm:erratum}(i), which reproduces all three displayed values
simultaneously from one derivation chain.

\section{The Frobenius-certificate theorem, formalized}\label{sec:cert}

The mathematical content, proven in \cite{CBcert}: for a prime $p$ with
$n < p \le 2n$, the Frobenius identity $k^p - k \equiv (k+j)^p - (k+j)
\pmod p$ makes $\varphi = (k^p - k)^2$ congruent to $(k+j)^2$ modulo
$(k+j)^6$ at every pole $-j$, while $\deg \varphi = 2p \le 4n + 2$ keeps
$\varphi R_n$ decaying at infinity; the residue-sum theorem on $\mathbb P^1$
then forces $\sum_j a_{3,j} \equiv 0 \pmod p$. In certificate form this is
the three-term relation with coefficients $(c_3, c_{p+2}, c_{2p+1}) =
(1, -2, 1)$ among the decay relations
\[
E_M \;=\; \sum_{i=1}^{\min(6,M)} \bin{-i}{M-i} \sum_{j} j^{M-i}\, a_{i,j}
\;=\; 0 \qquad (1 \le M \le 4n+4),
\]
which are the coefficients of the expansion of $R_n$ at infinity.

Conceptually: the certificate must prescribe a $6$-jet at each of the $n+1$
poles ($6(n+1)$ linear conditions) under a degree budget of $4n+3$ imposed by
the residue theorem --- an overdetermined Hermite interpolation problem in
characteristic zero, which is why the characteristic-zero $E_M$ relations
alone cannot force the congruence. Frobenius solves it in characteristic $p$
at degree $2p$, within budget precisely on the prime window; the mechanism is
\emph{Frobenius-assisted global jet interpolation plus residue duality}, and
it generalizes to arbitrary polynomial weights $Q$ with
$2p + \deg Q \le 4n+3$ \cite{CBcert}.

\begin{theorem}[{\lean{Cbcert/Main.lean}}]\label{thm:cert}
For every $n \ge 1$ and every prime $p$ with $5 \le p$ and $n < p \le 2n$,
verified by the Lean kernel with standard axioms:
\[
\text{\lean{res\_congruence\_w}}:\; p \mid w_n, \qquad
\text{\lean{res\_congruence\_wt}}:\; p \mid \wtd w_n, \qquad
\text{\lean{res\_congruence\_pn}}:\; p \mid p_n,
\]
where $p \mid x$ for a $p$-integral rational $x$ means $\mathrm{res}_p(x) =
0$ under the residue homomorphism $\mathrm{res}_p$ to $\Z/p$. Equivalently
\textup{(}\lean{pn\_valuation} etc.\textup{)}: $p_n = 0$ or
$\mathrm{ord}_p(p_n) \ge 1$.
\end{theorem}

The last statement of Theorem~\ref{thm:cert} is the clean-interval
central-binomial cancellation $(\mathrm{CB}_1)$: every prime of the window
$(n, 2n]$, $p \ge 5$, divides $p_n$.

\subsection*{Architecture}

The proof formalizes in four independent modules plus an assembly, mirroring
\cite{CBcert}:

\begin{center}
\begin{tabular}{lll}
module & content & status \\
\hline
\lean{Defs.lean} & all definitions (computable over $\Q$) & frozen, gated \\
\lean{Certificate.lean} & Lemma B: the $\Fp$ certificate identity & proven \\
\lean{PartialFraction.lean} & \lean{pf\_cleared}: the all-$n$ cleared identity & proven \\
\lean{Decay.lean} & Lemma A: $E_M(a) = E_M(\wtd a) = 0$ & proven \\
\lean{Integrality.lean} & Lemma C: $p$-integrality of $a, \wtd a, H$ & proven \\
\lean{Main.lean} & residue map, reordering, main theorems & proven \\
\lean{Numeric.lean} & kernel-reduction gates at $n = 2, 3$ & proven \\
\lean{ErrorExhibit.lean} & Theorem~\ref{thm:erratum} & proven \\
\end{tabular}
\end{center}

\emph{Lemma B} is pure binomial arithmetic in $\Z/p$: for $p \ge 5$,
$1 \le i \le 6$ and every $x \in \Z/p$,
\[
\sum_{M \in \{3,\, p+2,\, 2p+1\}} c_M \bin{-i}{M-i}\, x^{M-i}
\;=\; \delta_{i,3} \qquad\text{in } \Z/p,
\]
proven from $x^p = x$ and $p \mid \bin{p}{m}$ for $0 < m < p$. It is the
mathematically novel part of \cite{CBcert} and was the \emph{easiest} module
to formalize --- a data point for the certificate-shaped design thesis.

\emph{Lemma A} never needed the anticipated Laurent-series machinery: with
the cleared identity \lean{pf\_cleared} in hand, the decay relations become
the statement that $E_M(a)$ is the coefficient of $X^M$ in
$R_n(1/X) = X^{4n+5}\,\wtd N(X)/\wtd D(X)$ with $\wtd D(0) = 1$, an
elementary computation in $\Q[[X]]$. The load-bearing insight in
\lean{pf\_cleared} itself was that the truncated-multiplication primitive
used by the computable definitions \emph{is} the Cauchy product, so the
Taylor-coefficient definition of $a_{i,j}$ composes directly into the
polynomial identity.

\emph{Lemma C} bounds every denominator of $a_{i,j}$, $\wtd a_{i,j}$
($j \le n$) and $H_j^{(i)}$ by primes $\le n < p$: the pole differences
$|j - j'| \le n$ are the only sources.

\emph{Assembly.} A residue homomorphism $\mathrm{res}_p$ is developed on the
$p$-integral subring of $\Q$ (compatible with $+$, $\times$, and with
$\mathrm{res}_p(x) = 0 \iff x = 0 \vee \mathrm{ord}_p(x) \ge 1$), and a
coefficient-agnostic reordering lemma turns Lemmas A + B + C into
$\mathrm{res}_p\bigl(\sum_j b_{3,j}\bigr) = 0$ for both coefficient systems
$b \in \{a, \wtd a\}$ at once; $p \mid p_n$ then follows from
\eqref{eq:ladders} by ring arithmetic. All three certificate indices satisfy
$M \le 4n+4$ (respectively $4n+2$ for $\wtd a$) precisely because
$p \le 2n$; the case $p = 2n+1$ is genuinely excluded.

\section{What the formalization caught}\label{sec:caught}

Two statement-level errors in our own paper-level drafts were discovered and
repaired during formalization. We record them because both are invisible to
ordinary numerical testing on the intended domain, and both concern the
\emph{statements}, not the proofs.

\begin{enumerate}[label=\textup{(\arabic*)}]
\item \textbf{The valuation phrasing is unsound at zero.} The intended
statement $\mathrm{ord}_p(p_n) \ge 1$ was originally frozen as
$1 \le \lean{padicValRat}\ p\ (p_n)$. Mathlib defines
$\lean{padicValRat}\ p\ 0 = 0$, so this phrasing silently asserts
$p_n \ne 0$ --- a genuinely separate arithmetic fact (true numerically for
$n \le 60$; open in general, and \emph{not} needed for any downstream
denominator application). The canonical theorems are therefore stated in
residue form, with the faithful disjunction
$p_n = 0 \vee 1 \le \lean{padicValRat}\ p\ (p_n)$ provided as a corollary.
General nonvanishing is documented as an explicitly open problem, not a
hidden hypothesis.
\item \textbf{The integrality lemma needs its hypotheses.} As first staged,
Lemma C quantified over all pole indices $j$ and all primes $p > n$. Both
restrictions matter: the formal coefficient function evaluated out of range
or at $p = 2$ produces genuine counterexamples (at $n = 1$:
$a_{5,0} = 9/2$, and the out-of-range $a_{5,2} = 9/128$ with
$\mathrm{ord}_2 = -7$). The proven lemma carries $j \le n$ and $p$ odd; the
false general wrappers were deleted.
\end{enumerate}

Neither finding affects \cite{CBcert}, whose arguments always operate in
range with $p \ge 5$; both are exactly the kind of boundary defect that
survives refereeing.

\section{The second wave: finite range, the Lucas--Frobenius row, and the
conditional assembly}\label{sec:wave2}

Three further results were formalized on 19--20 July 2026, extending the
development from ``the band theorem plus the erratum'' toward the full
corrected law.

\subsection*{The corrected law for $n \le 8$, with nonvanishing}

Modules \lean{Cbcert/FiniteLaw/N2..N8.lean} verify, by pure kernel
computation (no \lean{native\_decide}), the full corrected law
$12\,d_n^5 P_n \in \Z$ for every $n \le 8$, together with $p_n \ne 0$ at each
such $n$ --- so on this range the disjunctive main theorems collapse to the
sharp valuation form unconditionally. An exact-arithmetic gate confirmed the
law out to $n = 24$ in Python before any Lean statement was generated, and
each Lean file checks a precomputed integer witness rather than searching.
(The kernel cost grows $\approx 1.75\times$ per level and does not
parallelize within a level; $n \le 8$ was chosen deliberately, and the
independent adversarial sweep described below now covers $n \le 360$.)

\subsection*{The Lucas--Frobenius row}

For the manifestly integral Brown--Zudilin double-binomial coefficient
\[
Q_n \;=\; \sum_{k,l=0}^{n} \bin{n+k}{n}\bin nk^2 \bin{n+l}{n}\bin nl^2
\bin{n+k+l}{n},
\]
module \lean{Cbcert/LucasRow.lean} proves, sorry-free with standard axioms:

\begin{quote}\small\ttfamily
theorem Q\_lucas\_frobenius (a r p : $\mathbb N$) (hp : p.Prime) (hr : r < p) :\\
\hspace*{1em}(Q (a*p + r) : $\Z$) $\equiv$ (Q a : $\Z$) * (Q r : $\Z$) [ZMOD p]
\end{quote}

\noindent
This is the ``$q$-row'' of the conjectural Frobenius structure on the period
lattice: the diagonal of the descent, proven unconditionally. The proof uses
Mathlib's Lucas theorem
(\lean{Choose.choose\_modEq\_choose\_mod\_mul\_choose\_div\_nat}) per binomial
factor. Formalization again caught a genuine gap in the human proof sketch:
on summands with a base-$p$ carry, the five per-factor Lucas splits do
\emph{not} hold simultaneously, so one cannot simply multiply five
congruences; the correct argument exhibits a single common vanishing factor
on the carry branch. The sketch was right in spirit and wrong in mechanism
--- precisely the kind of defect that survives refereeing.

\subsection*{The conditional assembly: one hypothesis from the full law}

Module \lean{Cbcert/Assembly.lean} formalizes the master reduction: define

\begin{quote}\small\ttfamily
def DescentHyp : Prop :=\\
\hspace*{1em}$\forall$ (m p : $\mathbb N$), p.Prime $\to$ 5 $\le$ p $\to$ p $\le$ m $\to$\\
\hspace*{2em}padicValRat p (P (m / p)) $-$ 5 $\le$ padicValRat p (P m)
\end{quote}

\noindent
--- exactly the descent $(\mathrm D)$: one base-$p$ digit down costs at most
five $p$-adic orders. Then, sorry-free with standard axioms:

\begin{quote}\small\ttfamily
theorem descent\_law (hD : DescentHyp) :\\
\hspace*{1em}$\forall$ n $\ge$ 1, $\forall$ p prime, 5 $\le$ p $\to$
($-$5)$\cdot$(Nat.log p n) $\le$ padicValRat p (P n)
\end{quote}

\noindent
The proof iterates $(\mathrm D)$ down the digit chain and discharges both
terminal cases from theorems already in the development: $2n_L < p$ via the
integrality lemmas (Lemma C), and $n_L < p \le 2n_L$ via the certificate
theorem (Theorem~\ref{thm:cert}) plus a one-carry Kummer lemma. Two design
points matter. First, the statement is \emph{trap-free}: under Mathlib's
$\lean{padicValRat}\ p\ 0 = 0$ convention the induction closes with no
nonvanishing hypotheses anywhere. Second, \lean{hD} is a \emph{hypothesis},
not an axiom --- the axiom audit shows the standard triple only, and the
theorem states formally and exactly what remains open: prove $(\mathrm D)$,
and the full corrected law at every $p \ge 5$ follows by composition of
kernel-verified pieces.

$(\mathrm D)$ itself is the object of an active proof campaign (with the
collaborating model Sol, a GPT-5.6 instance, whose sustained multi-session
work produced the entire Phase~2 reduction summarized here and in the
companion progress note), which has reduced it, on paper, through a chain of
exact certificate-shaped reductions --- per-digit gates at the true
singularities of the underlying Ap\'ery-like recurrence, then a proved $a=1$
midpoint gate, with the remaining difficulty localized to a single period
connection coefficient --- each step verified in exact arithmetic before use. The adversarial evidence base for
$(\mathrm D)$ now stands at $n \le 360$, $p \le 73$: $5{,}989$ exact
descents, zero violations, including the previously untested three-digit
descent ranges, with the law riding the boundary (slack $0$) at $p = 2$ for
every $n$ --- the constant $12$ is exactly sharp throughout. Finite evidence,
not proof; but the formal statement of what it is evidence \emph{for} is now
frozen in \lean{DescentHyp}.

\section{Scope and open problems}\label{sec:scope}

The formalized results are: the erratum and its sharp $n = 2$ repair
(Theorem~\ref{thm:erratum}); the clean-interval theorem
(Theorem~\ref{thm:cert}); the finite-range law with nonvanishing for
$n \le 8$; the Lucas--Frobenius row; and the conditional assembly
$(\mathrm D) \Rightarrow$ full law at $p \ge 5$
(Section~\ref{sec:wave2}). Explicitly \emph{not} formalized, because they
are not yet theorems on paper: the descent $(\mathrm D)$ itself (the single
remaining gap, object of the active campaign described above); the
$\{2,3\}$-part of the constant $12$ (the fixed losses at $p = 2, 3$, a
separate finite obligation); and the general nonvanishing
$w_n, \wtd w_n, p_n \ne 0$ (not needed by anything above).

\section{Verification}\label{sec:verify}

The development lives in the \lean{cbcert/} directory of the public
repository \cite{repo}, pinned to Lean toolchain
\lean{v4.33.0-rc1} and a fixed Mathlib commit. To verify:
\lean{lake exe cache get \&\& lake build} (the full build reports
$\approx 8680$ jobs and finishes with no \lean{sorry} warnings), then
\lean{\#print axioms} on any theorem named above — including
\lean{Q\_lucas\_frobenius}, \lean{descent\_law}, and the \lean{FiniteLaw}
aggregate \lean{law\_upto}; each reports exactly
\lean{[propext, Classical.choice, Quot.sound]} (in particular no
\lean{sorryAx} in \lean{descent\_law}: the descent enters as a named
hypothesis, not an axiom). The repository README
contains the recipe, including which displays of \cite{BZ22} to compare
against which Lean docstrings.

\subsection*{Provenance}

The formalization was produced on 19 July 2026 in a single orchestrated
session by a fleet of LLM agents (Claude, Anthropic), following the
statement-first protocol described above: definitions frozen against exact
numeric gates before lemma work, statements corrected only under the
documented-erratum discipline of Section~\ref{sec:caught}. As with any
formal development, correctness rests on the Lean kernel and the pinned
toolchain, not on the process that produced the proof scripts.

\begin{thebibliography}{9}

\bibitem{BZ22}
\textsc{F.~Brown, W.~Zudilin},
On cellular rational approximations to $\zeta(5)$,
arXiv:2210.03391 [math.NT] (2022; revised 26 January 2026).

\bibitem{Zu02}
\textsc{W.~Zudilin},
A third-order Ap\'ery-like recursion for $\zeta(5)$,
\emph{Mat. Zametki} \textbf{72} (2002), no.~5, 796--800;
English transl., \emph{Math. Notes} \textbf{72} (2002), no.~5--6, 733--737.

\bibitem{CBcert}
A Frobenius certificate for central-binomial cancellation in Zudilin's
symmetric $\zeta(5)$ construction, companion paper, this repository
(\lean{worthiness/cb\_certificate.tex}), 17 July 2026; revised 19 July 2026.

\bibitem{Lean4}
\textsc{L.~de Moura, S.~Ullrich},
The Lean 4 theorem prover and programming language,
\emph{CADE 28}, Lecture Notes in Comput. Sci. \textbf{12699} (2021), 625--635.

\bibitem{mathlib}
\textsc{The mathlib Community},
The Lean mathematical library,
\emph{CPP 2020}, ACM (2020), 367--381.

\bibitem{repo}
Odd zeta values: machine verifications (repository),
\url{https://github.com/rain-1/odd-zeta-values-autoformalization}.

\end{thebibliography}

\end{document}
